\textbf{Phần 1. Trắc nghiệm lựa chọn}
\setcounterpageref{ex}{0}

\begin{ex} %[5H5N27]
Công thức tính chu vi $C$ của hình tròn có bán kính $r$ là
\choice
{$C = 3{,}14 \times r$}
{$C = 3{,}14 \times r \times r$}
{$C = 3{,}14 \times 4 \times r$}
{\True $C = 3{,}14 \times 2 \times r$}
\loigiai{
Chu vi hình tròn bằng tích của số $3{,}14$ với $2$ lần bán kính: $C = 3{,}14 \times 2 \times r$.\\
\textbf{Đáp án của câu hỏi là D}
}
\end{ex}

\begin{ex} %[5H5N27]
Công thức tính diện tích $S$ của hình tròn có bán kính $r$ là
\choice
{\True $S = 3{,}14 \times r \times r$}
{$S = 3{,}14 \times 2 \times r$}
{$S = 3{,}14 \times d$}
{$S = 3{,}14 \times 2 \times r \times r$}
\loigiai{
Diện tích hình tròn bằng số $3{,}14$ nhân với bán kính rồi nhân với bán kính: $S = 3{,}14 \times r \times r$.\\
\textbf{Đáp án của câu hỏi là A}
}
\end{ex}

\begin{ex} %[5H5H27]
Chu vi của hình tròn có đường kính $6\text{ dm}$ là
\choice
{$18\text{ dm}$}
{$37{,}68\text{ dm}$}
{$28{,}26\text{ dm}$}
{\True $18{,}84\text{ dm}$}
\loigiai{
Chu vi hình tròn là: $3{,}14 \times 6 = 18{,}84\text{ (dm)}$.\\
\textbf{Đáp án của câu hỏi là D}
}
\end{ex}

\begin{ex} %[5H5H27]
Diện tích của hình tròn có bán kính $10\text{ cm}$ là
\choice
{$31{,}4\text{ cm}^2$}
{\True $314\text{ cm}^2$}
{$62{,}8\text{ cm}^2$}
{$3\,140\text{ cm}^2$}
\loigiai{
Diện tích hình tròn là: $3{,}14 \times 10 \times 10 = 314\text{ (cm}^2\text{)}$.\\
\textbf{Đáp án của câu hỏi là B}
}
\end{ex}

\begin{ex} %[5H5H27]
Diện tích của một tấm thảm hình tròn có đường kính $20\text{ dm}$ là
\choice
{$62{,}8\text{ dm}^2$}
{$1\,256\text{ dm}^2$}
{\True $314\text{ dm}^2$}
{$31{,}4\text{ dm}^2$}
\loigiai{
Bán kính tấm thảm là: $20 : 2 = 10\text{ (dm)}$.\\
Diện tích tấm thảm là: $3{,}14 \times 10 \times 10 = 314\text{ (dm}^2\text{)}$.\\
\textbf{Đáp án của câu hỏi là C}
}
\end{ex}

\begin{ex} %[5H5H27]
Bãi thả khí cầu của một công ty có dạng hình tròn bán kính $200\text{ m}$. Diện tích của bãi thả khí cầu đó là
\choice
{\True $125\,600\text{ m}^2$}
{$1\,256\text{ m}^2$}
{$12\,560\text{ m}^2$}
{$1\,256\,000\text{ m}^2$}
\loigiai{
Diện tích bãi thả khí cầu là: $3{,}14 \times 200 \times 200 = 125\,600\text{ (m}^2\text{)}$.\\
\textbf{Đáp án của câu hỏi là A}
}
\end{ex}

\begin{ex} %[5H5H27]
Một hình tròn có chu vi là $31{,}4\text{ cm}$. Bán kính của hình tròn đó là
\choice
{$10\text{ cm}$}
{$2{,}5\text{ cm}$}
{\True $5\text{ cm}$}
{$15{,}7\text{ cm}$}
\loigiai{
Bán kính hình tròn là: $31{,}4 : 3{,}14 : 2 = 5\text{ (cm)}$.\\
\textbf{Đáp án của câu hỏi là C}
}
\end{ex}

\begin{ex} %[5H5V27]
Cái cạp rổ hình tròn có đường kính $50\text{ cm}$. Chiều dài phần hai đầu tre buộc vào nhau là $15\text{ cm}$. Chiều dài của thanh tre uốn thành cái cạp rổ là
\choice
{\True $172\text{ cm}$}
{$157\text{ cm}$}
{$182\text{ cm}$}
{$329\text{ cm}$}
\loigiai{
Chu vi cạp rổ là: $3{,}14 \times 50 = 157\text{ (cm)}$.\\
Chiều dài thanh tre là: $157 + 15 = 172\text{ (cm)}$.\\
\textbf{Đáp án của câu hỏi là A}
}
\end{ex}

\begin{ex} %[5H5V27]
Một sợi dây thừng quấn quanh một gốc cây hình tròn có đường kính $4\text{ dm}$ đúng $3$ vòng. Phần dây không quấn vào thân cây dài $2{,}8\text{ m}$. Chiều dài của sợi dây thừng đó là
\choice
{$3{,}768\text{ m}$}
{\True $6{,}568\text{ m}$}
{$40{,}48\text{ m}$}
{$65{,}68\text{ m}$}
\loigiai{
Chu vi gốc cây là: $3{,}14 \times 4 = 12{,}56\text{ (dm)} = 1{,}256\text{ (m)}$.\\
Chiều dài sợi dây thừng là: $1{,}256 \times 3 + 2{,}8 = 3{,}768 + 2{,}8 = 6{,}568\text{ (m)}$.\\
\textbf{Đáp án của câu hỏi là B}
}
\end{ex}

\begin{ex} %[5H5V27]
Ở một vùng sa mạc, người ta trồng lúa trên những thửa ruộng hình tròn bán kính $50\text{ m}$. Biết có $1\,000$ thửa ruộng như vậy. Tổng diện tích trồng lúa là
\choice
{$78\,500\text{ m}^2$}
{$785\,000\text{ m}^2$}
{$31\,400\,000\text{ m}^2$}
{\True $7\,850\,000\text{ m}^2$}
\loigiai{
Diện tích $1$ thửa ruộng là: $3{,}14 \times 50 \times 50 = 7\,850\text{ (m}^2\text{)}$.\\
Tổng diện tích $1\,000$ thửa ruộng là: $7\,850 \times 1\,000 = 7\,850\,000\text{ (m}^2\text{)}$.\\
\textbf{Đáp án của câu hỏi là D}
}
\end{ex}

\begin{ex} %[5H5V27]
Một sân vận động được tạo bởi một hình chữ nhật có kích thước $100\text{ m} \times 70\text{ m}$ và hai nửa hình tròn có đường kính $70\text{ m}$ ở hai đầu. Diện tích của sân vận động đó là
\choice
{$7\,000\text{ m}^2$}
{$14\,693\text{ m}^2$}
{\True $10\,846{,}5\text{ m}^2$}
{$10\,146{,}5\text{ m}^2$}
\loigiai{
Diện tích phần hình chữ nhật là: $100 \times 70 = 7\,000\text{ (m}^2\text{)}$.\\
Bán kính hai nửa hình tròn là: $70 : 2 = 35\text{ (m)}$.\\
Diện tích hai nửa hình tròn (tạo thành một hình tròn) là: $3{,}14 \times 35 \times 35 = 3\,846{,}5\text{ (m}^2\text{)}$.\\
Diện tích sân vận động là: $7\,000 + 3\,846{,}5 = 10\,846{,}5\text{ (m}^2\text{)}$.\\
\textbf{Đáp án của câu hỏi là C}
}
\end{ex}

\begin{ex} %[5H5V27]
Một giếng nước có miệng giếng là hình tròn bán kính $8\text{ dm}$. Người ta xây thành giếng rộng $3\text{ dm}$ bao quanh miệng giếng. Diện tích phần đất xây thành giếng đó là
\choice
{$200{,}96\text{ dm}^2$}
{\True $178{,}98\text{ dm}^2$}
{$379{,}94\text{ dm}^2$}
{$150{,}72\text{ dm}^2$}
\loigiai{
Bán kính hình tròn lớn (bao gồm cả thành giếng) là: $8 + 3 = 11\text{ (dm)}$.\\
Diện tích hình tròn lớn là: $3{,}14 \times 11 \times 11 = 379{,}94\text{ (dm}^2\text{)}$.\\
Diện tích miệng giếng là: $3{,}14 \times 8 \times 8 = 200{,}96\text{ (dm}^2\text{)}$.\\
Diện tích thành giếng là: $379{,}94 - 200{,}96 = 178{,}98\text{ (dm}^2\text{)}$.\\
\textbf{Đáp án của câu hỏi là B}
}
\end{ex}

\textbf{Phần 2. Trắc nghiệm đúng sai}
\setcounterpageref{ex}{0}

\begin{ex} %[5H5H27]
Xét tính đúng sai của các khẳng định sau:
\choiceTF
{\True Khi bán kính hình tròn tăng lên gấp đôi thì chu vi hình tròn tăng lên gấp đôi.}
{Khi bán kính hình tròn tăng lên gấp đôi thì diện tích hình tròn tăng lên gấp đôi.}
{\True Chu vi của hình tròn có bán kính $5\text{ cm}$ là $31{,}4\text{ cm}$.}
{\True Diện tích của hình tròn có bán kính $3\text{ dm}$ là $28{,}26\text{ dm}^2$.}
\loigiai{
\begin{itemize}
    \item $C = 3{,}14 \times 2 \times r \Rightarrow$ tỉ lệ thuận bậc 1 với $r$, nên $r$ gấp 2 thì chu vi gấp 2.
    \item $S = 3{,}14 \times r \times r \Rightarrow$ khi $r$ gấp 2 thì $S$ gấp $2 \times 2 = 4$ lần (khẳng định tăng 2 lần là sai).
    \item $C = 3{,}14 \times 2 \times 5 = 31{,}4\text{ cm}$.
    \item $S = 3{,}14 \times 3 \times 3 = 28{,}26\text{ dm}^2$.
\end{itemize}
\begin{itemize}
    \item a) \textbf{Đúng}.
    \item b) \textbf{Sai}.
    \item c) \textbf{Đúng}.
    \item d) \textbf{Đúng}.
\end{itemize}
}
\end{ex}

\begin{ex} %[5H5V27]
Chú rùa vàng bò từ $A$ đến $B$ theo đường nửa đường tròn lớn màu đỏ (đường kính $9\text{ dm}$), chú rùa nâu bò từ $B$ đến $A$ theo $3$ nửa đường tròn nhỏ màu xanh (có đường kính lần lượt là $3\text{ dm}, 1\text{ dm}, 5\text{ dm}$). Xét tính đúng sai của các khẳng định sau:
\choiceTF
{\True Quãng đường rùa vàng đi là $14{,}13\text{ dm}$.}
{\True Quãng đường rùa nâu đi là $14{,}13\text{ dm}$.}
{Rùa vàng đi quãng đường dài hơn rùa nâu.}
{\True Quãng đường đi của hai chú rùa là bằng nhau.}
\loigiai{
\begin{itemize}
    \item Nửa chu vi đường tròn lớn: $3{,}14 \times 9 : 2 = 14{,}13\text{ dm}$.
    \item Tổng 3 nửa chu vi nhỏ: $3{,}14 \times (3 + 1 + 5) : 2 = 3{,}14 \times 9 : 2 = 14{,}13\text{ dm}$.
    \item Quãng đường hai chú rùa đi bằng nhau.
\end{itemize}
\begin{itemize}
    \item a) \textbf{Đúng}.
    \item b) \textbf{Đúng}.
    \item c) \textbf{Sai}.
    \item d) \textbf{Đúng}.
\end{itemize}
}
\end{ex}

\textbf{Phần 3. Trả lời ngắn}
\setcounterpageref{ex}{0}

\begin{ex} %[5H5H27]
Tính chu vi của hình tròn có đường kính $10\text{ m}$ theo mét. (Nhập kết quả dưới dạng số thập phân).
\shortans{31,4}
\loigiai{
\textbf{Đáp số:} $31{,}4$\\
$3{,}14 \times 10 = 31{,}4\text{ (m)}$.
}
\end{ex}

\begin{ex} %[5H5H27]
Tính diện tích của hình tròn có bán kính $2\text{ cm}$ theo xăng-ti-mét vuông. (Nhập kết quả dưới dạng số thập phân).
\shortans{12,56}
\loigiai{
\textbf{Đáp số:} $12{,}56$\\
$3{,}14 \times 2 \times 2 = 12{,}56\text{ (cm}^2\text{)}$.
}
\end{ex}

\begin{ex} %[5H5H27]
Một hình tròn có chu vi là $62{,}8\text{ cm}$. Tính diện tích của hình tròn đó theo xăng-ti-mét vuông.
\shortans{314}
\loigiai{
\textbf{Đáp số:} $314$\\
Bán kính: $62{,}8 : 3{,}14 : 2 = 10\text{ cm}$. Diện tích: $3{,}14 \times 10 \times 10 = 314\text{ cm}^2$.
}
\end{ex}

\begin{ex} %[5H5V27]
Một sợi dây thừng quấn quanh gốc cây hình tròn đường kính $4\text{ dm}$ đủ $3$ vòng. Phần thừa dài $2{,}8\text{ m}$. Hỏi cả sợi dây dài bao nhiêu mét? (Nhập kết quả dưới dạng số thập phân).
\shortans{6,568}
\loigiai{
\textbf{Đáp số:} $6{,}568$\\
Chu vi 1 vòng: $3{,}14 \times 0{,}4 = 1{,}256\text{ m}$. Dây dài: $1{,}256 \times 3 + 2{,}8 = 6{,}568\text{ m}$.
}
\end{ex}

\begin{ex} %[5H5V27]
Sân vận động gồm hình chữ nhật $100\text{ m} \times 70\text{ m}$ và hai nửa hình tròn đường kính $70\text{ m}$. Tính diện tích sân vận động đó theo mét vuông. (Nhập kết quả dưới dạng số thập phân).
\shortans{10846,5}
\loigiai{
\textbf{Đáp số:} $10846{,}5$\\
$100 \times 70 + 3{,}14 \times 35 \times 35 = 7\,000 + 3\,846{,}5 = 10\,846{,}5\text{ (m}^2\text{)}$.
}
\end{ex}

\begin{ex} %[5H5V27]
Miệng giếng hình tròn bán kính $8\text{ dm}$, thành giếng bao quanh rộng $3\text{ dm}$. Tính diện tích phần đất xây thành giếng theo đề-xi-mét vuông. (Nhập kết quả dưới dạng số thập phân).
\shortans{178,98}
\loigiai{
\textbf{Đáp số:} $178{,}98$\\
$3{,}14 \times (11^2 - 8^2) = 3{,}14 \times 57 = 178{,}98\text{ (dm}^2\text{)}$.
}
\end{ex}
